% !PRINTED 418
\weberpage{418}
%<<CONT
\gesperrt{polyhedral groups} must be identical. We wish here to treat in more detail only the most interesting case, that $P$ is the \gesperrt{icosahedral group}, which moreover we may assume, without restriction of generality, in the normal form (\S.\,58).

For this group we have derived the three fundamental forms in \S.\,60. If we denote the variables of these fundamental forms by $y_1, y_2$, they are the three homogeneous functions of the $12$th, $20$th and $30$th degrees:
\weq{9}{\begin{aligned}
f(y) &= y_1 y_2 (y_1^{10} + 11 y_1^5 y_2^5 - y_2^{10}) \\
H(y) &= -y_1^{20} - y_2^{20} + 228 (y_1^{15} y_2^5 - y_2^{15} y_1^5) - 494 y_1^{10} y_2^{10} \\
T(y) &= y_1^{30} + y_2^{30} + 522 (y_1^{25} y_2^5 - y_1^5 y_2^{25}) - 10005 (y_1^{20} y_2^{10} + y_2^{20} y_1^{10}),
\end{aligned}}
between which the relation
\weq{10}{T^2 + H^3 = 1728 f^5}
further subsists.

If in the quotient $T^2 : f^5$ we put for the ratio $y_1 : y_2$ one of the 60 values $\vartheta_\alpha$ that arise from $\vartheta$ by the icosahedral substitutions, this quotient always receives the same value, which can therefore be expressed rationally in terms of the coefficients of equation (8), that is, rationally in terms of $z$. We may set it directly equal to $z$, and obtain for the resolvent (8) the form
\weq{11}{T^2 - zf^5 = 0,}
which on account of (10) is equivalent to
\[
H^3 - (1728 - z) f^5 = 0
\]
This equation is homogeneous in $y_1, y_2$ and of the $60$th degree, and if we set $y_2 = 1$, its 60 roots are $y_1 = \vartheta_\alpha$. It is the \gesperrt{icosahedral equation}, which was already defined in \S.\,60, and which therefore, written out in full in the form (11), runs thus:
\[
(\vartheta^{30} + 522 \vartheta^{25} - 10005 \vartheta^{20} - 10005 \vartheta^{10} - 522 \vartheta^5 + 1)^2 = z \vartheta^5 (\vartheta^{10} + 11 \vartheta^5 - 1)^5.
\]

The problem, as we have now framed it, comes down to the question:

\gesperrt{Which equations can be solved by means of the icosahedral equation?}
